% \chapter{Concluzii}
% %\chapter*{Concluzii}
% \label{conclusions}

% \par Concluzii ...
\chapter{Conclusions}
\label{ch:conclusions}

This thesis investigated how a commodity iOS device can support \textbf{interior design preview} through \textbf{parametric room capture}, \textbf{mesh-based editing}, and optional \textbf{plenoptic visualisation}. The theoretical chapters argued that Apple's Object Capture instantiates a classical \textbf{photogrammetric} chain (features, SfM, MVS-class dense matching, Poisson-class surface reconstruction, texturing), whereas RoomPlan builds on \textbf{depth fusion} and learning to emit a structured \texttt{CapturedRoom}. The second theory chapter positioned \textbf{3D Gaussian Splatting} \cite{kerbl20233dgs} as a real-time differentiable alternative to meshes for \emph{viewing}, not for CAD-style furniture replacement.

\section{Summary of Results}

The software artefact demonstrates a coherent pipeline: \textbf{RoomPlan} $\rightarrow$ serialised scan $\rightarrow$ \textbf{RealityKit} editor $\rightarrow$ USDZ replacements with gestures $\rightarrow$ persistence across sessions. \textbf{Object Capture} extends asset authoring; a \textbf{Metal} splat renderer implements R-7 for imported \texttt{.ply} data. Requirements R-1--R-7 (and associated non-functional goals) are addressed at the design level in Chapter~\ref{ch:requirements} and at the code level in Chapter~\ref{ch:implementation}.

\section{Limitations}

Accuracy is bounded by RoomPlan detection quality, photogrammetry assumptions in Object Capture, and the lack of closed-loop feedback between splats and parametric room semantics. 3DGS training remains external; real-time splat \emph{training} on device is research-grade \cite{Guo26}. Usability has not been quantified with SUS/HARUS in this thesis text; Kandil et al.\ \cite{kandil2021ux} provide a methodology for a future study.

\section{Future Work}

Promising extensions include: formal user evaluation; tighter integration between exported scan images and external COLMAP/3DGS pipelines; level-of-detail for large splats; semantic suggestions informed by work such as K\'an et al.\ \cite{kan2021ar} Mithun et al. \cite{Geo2023} demonstrate that synthetic datasets with rich semantic labels can substantially improve indoor scene understanding for AR applications, a direction complementary to the real-scan approach taken here.

The project shows that \textbf{merging API-level room semantics with an editable mesh workflow} yields a practical middle ground between live AR stickers and professional BIM tools, suitable for iterative interior exploration on mobile hardware.
